\begin{tabular}{@{}rlp{0.5\textwidth}@{}}
\toprule
\thead{°C / °F} & \thead{Phase}           & \thead{Beschreibung} \\ \midrule
60 / 140        & Sterilisation           & Die Temperatur ist zu hoch für Ihre Mikroorganismen und sie sterben ab.\\
75 / 167        & Gelbildung            & Ein Gel bildet sich auf der Oberfläche und bewahrt die Teigstruktur.
                                            Er ist noch dehnbar und kann im Ofen aufgehen.\\
100 / 212       & Wasserverdampfung       & Wasser beginnt zu verdampfen und bläht die Alveolen des Teigs auf.\\
118 / 244       & Essigsäureverdampfung & Die essigartig schmeckende Säure beginnt zu verdampfen, die Säure nimmt ab.\\
122 / 252       & Milchsäureverdampfung & Die milchig schmeckende Milchsäure beginnt zu verdampfen, die Säure nimmt weiter ab.\\
140 / 284       & Maillard-Reaktion       & Die Maillard-Reaktion beginnt Stärken und Proteine zu verändern.
                                            Der Teig beginnt zu bräunen.\\
170 / 338       & Karamellisierung          & Verbleibende Zucker beginnen zu karamellisieren und verleihen Ihrem Brot einen besonderen Geschmack.\\ \bottomrule
\end{tabular}
